\begin{center}\begin{tikzpicture}[font=\small,>=stealth]
\draw[black!30,rounded corners=2pt] (0,0) rectangle (14.4,8.2);
\node[anchor=west,font=\small\bfseries] at (.5,7.72) {例 3：绝对值方程的两个正根，还要检查几何存在性};
\node at (3.6,6.65) {$PF_1=7,\ PF_2=1$};\begin{scope}[shift={(3.6,3.8)},scale=.145]
\draw[orange!85!black,thick] (-5,0) circle (7);
\draw[violet,thick] (5,0) circle (1);
\node[fill=black,circle,inner sep=1.5pt] at (-5,0) {};\node[below=4pt] at (-5,0) {$F_1$};
\node[fill=black,circle,inner sep=1.5pt] at (5,0) {};\node[below=4pt] at (5,0) {$F_2$};
\end{scope}
\node at (3.6,1.3) {两圆分离：不存在这样的点};
\node at (10.8,6.65) {$PF_1=7,\ PF_2=13$};\begin{scope}[shift={(10.8,3.8)},scale=.145]
\draw[orange!85!black,thick] (-5,0) circle (7);
\draw[violet,thick] (5,0) circle (13);
\node[fill=black,circle,inner sep=1.5pt] at (-5,0) {};\node[below=4pt] at (-5,0) {$F_1$};
\node[fill=black,circle,inner sep=1.5pt] at (5,0) {};\node[below=4pt] at (5,0) {$F_2$};
\node[fill=black,circle,inner sep=2pt] at (-6,6.9282032303) {};\node[above left=5pt] at (-6,6.9282032303) {$P$};
\end{scope}
\node at (10.8,1.3) {两圆相交：候选距离可实现};

\node[anchor=west] at (.7,.55) {同尺度；焦距 $10$。左图 $7+1<10$；右图交点同时满足两条距离要求};
\end{tikzpicture}\end{center}
